\documentclass[11pt]{article}
\usepackage[T1]{fontenc}
\usepackage[utf8]{inputenc}
\usepackage[landscape,margin=0.75in]{geometry}
\usepackage{lmodern}
\usepackage{microtype}
\usepackage{amsmath}
\usepackage{amssymb}
\usepackage{amsthm}
\usepackage{booktabs}
\usepackage{longtable}
\usepackage{array}
\usepackage{hyperref}
\usepackage{xurl}
\usepackage{orcidlink}
\setlength{\emergencystretch}{3em}

\renewcommand{\texttt}[1]{{\ttfamily\hyphenchar\font=`\-\relax
   \def\_{\textunderscore\penalty100 }#1}}

\title{Six Birds: No-Go Theorems for Audited Emergence: Supplement}
\author{Ioannis Tsiokos\,\orcidlink{0009-0009-7659-5964}}
\date{March 23, 2026}

\begin{document}
\maketitle
\section*{How to use this supplement}
This supplement collects the repo-facing apparatus that accompanies the main theorem paper.
It has two purposes.
First, it records provenance, evidence bindings, reproducibility commands, and fragility guardrails in one place.
Second, it summarizes the distribution of Lean support across the eight theorem fronts without interrupting the mathematical flow of the main text.

The main paper is intended to stand on its own as a theorem paper.
The present supplement is therefore organized around audit and provenance questions: where each claim is bound, which frozen assets support it, which reproduce steps build the relevant summaries, and where the formal boundary is explicitly frozen.
 \section{Provenance narrative}
The theorem line is assembled from a fixed project spine.
The theorem order and section homes come from \texttt{docs/project/theorem\_atlas.yaml}.
The formal statements used in the LaTeX theorem blocks come from \texttt{docs/project/theorem\_statement\_sheet.yaml}, and notation is stabilized in \texttt{docs/project/notation\_registry.yaml}.

Claim-level bindings are frozen in \texttt{docs/project/claim\_audit\_freeze.yaml} and \texttt{docs/project/theorem\_evidence\_map.yaml}.
Asset choices are frozen in \texttt{docs/project/paper\_asset\_freeze.yaml} and \texttt{docs/project/caption\_registry.yaml}.
Readiness and formal-support status are tracked in \texttt{docs/project/readiness\_checklist.yaml} together with the per-front Lean bridge files under \texttt{docs/project/lean\_*.yaml}.

The supplement tables below therefore do not introduce new mathematical claims.
They record how the frozen theorem fronts, claim rows, assets, and executable audit artifacts are aligned in the packaged repo state.
 \section{Theorem--evidence map}
This section records the frozen claim bindings, evidence handles, and asset/callout decisions used by the theorem fronts.
The first table is generated from \texttt{docs/project/theorem\_evidence\_map.yaml}, \texttt{docs/project/paper\_asset\_freeze.yaml}, and \texttt{docs/project/caption\_registry.yaml}.
The second table summarizes the frozen callout bindings and their primary sources.

{\footnotesize
\begin{longtable}{>{\raggedright\arraybackslash}p{0.13\textwidth} >{\raggedright\arraybackslash}p{0.17\textwidth} >{\raggedright\arraybackslash}p{0.18\textwidth} >{\raggedright\arraybackslash}p{0.24\textwidth} >{\raggedright\arraybackslash}p{0.20\textwidth}}
\caption{Claim-level theorem--evidence bindings.}\label{tab:evidence-map}\\
\toprule
Theorem & Claim & Primary handle & Evidence artifacts & Asset / callout decision\\
\midrule
\endfirsthead
\toprule
Theorem & Claim & Primary handle & Evidence artifacts & Asset / callout decision\\
\midrule
\endhead
\midrule
\multicolumn{5}{r}{Continued on next page}\\
\midrule
\endfoot
\bottomrule
\endlastfoot
\texttt{NG\_ARROW\_DPI} & \texttt{NG\_ARROW\_DPI.core} & \texttt{results/T08/arrow\_metrics.json} & \texttt{results/T08/arrow\_metrics.json}; \texttt{results/T15/initial\_distribution\_sweep.csv} & FIG-2:repo\_asset; Figure 2 ([FIG-2]). Hidden-clock DPI / protocol matrix: honest observ…\\
\texttt{NG\_PROTOCOL\_TRAP} & \texttt{NG\_PROTOCOL\_TRAP.core} & \texttt{results/T15/initial\_distribution\_sweep.csv} & \texttt{results/T15/initial\_distribution\_sweep.csv}; \texttt{results/T08/arrow\_metrics.json} & FIG-2:repo\_asset; Figure 2 ([FIG-2]). Hidden-clock DPI / protocol matrix: honest observ…\\
\texttt{NG\_FORCE\_FOREST} & \texttt{NG\_FORCE\_FOREST.core} & \texttt{results/T14/primitive\_coverage.json} & \texttt{results/T14/primitive\_coverage.json}; \texttt{results/master/witness\_metrics.csv} & FIG-3:repo\_asset; Figure 3 ([FIG-3]). Graph / affinity contrasts: forest support remove…\\
\texttt{NG\_FORCE\_NULL} & \texttt{NG\_FORCE\_NULL.core} & \texttt{results/master/witness\_metrics.csv} & \texttt{results/master/witness\_metrics.csv} & FIG-3:repo\_asset; Figure 3 ([FIG-3]). Graph / affinity contrasts: forest support remove…\\
\texttt{NG\_MACRO\_CLOSURE\_DEFICIT} & \texttt{NG\_MACRO\_CLOSURE\_DEFICIT.core} & \texttt{results/T09/closure\_metrics.json} & \texttt{results/T09/closure\_metrics.json}; \texttt{results/T09/grid\_check.json} & FIG-4:repo\_asset; Figure 4 ([FIG-4]). Closure front: exact packaged future laws define …\\
\texttt{NG\_OBJECT\_CONTRACTIVE} & \texttt{NG\_OBJECT\_CONTRACTIVE.core} & \texttt{results/T10/objecthood\_metrics.json} & \texttt{results/T10/objecthood\_metrics.json}; \texttt{results/T10/stability\_grid.json} & FIG-5:repo\_asset; Figure 5 ([FIG-5]). Objecthood front: in the admissible total-variati…\\
\texttt{NG\_LADDER\_IDEM} & \texttt{NG\_LADDER\_IDEM.core} & \texttt{results/T11/no\_ladder\_demos.json} & \texttt{results/T11/no\_ladder\_demos.json}; \texttt{results/T11/definability\_metrics.json} & FIG-6:repo\_asset; Figure 6 ([FIG-6]). Ladder front: fixed idempotent packaging saturate…\\
\texttt{NG\_LADDER\_BOUNDED\_INTERFACE} & \texttt{NG\_LADDER\_BOUNDED\_INTERFACE.core} & \texttt{results/T11/definability\_metrics.json} & \texttt{results/T11/definability\_metrics.json}; \texttt{results/T11/no\_ladder\_demos.json} & TBL-6:repo\_asset; Table 6 ([TBL-6]). Ladder theorem evidence, count core, and corollary…\\
\texttt{NG\_ARROW\_DPI} & \texttt{NG\_ARROW\_DPI.guardrail} & \texttt{results/T15/adversarial\_cases.json} & \texttt{results/T15/adversarial\_cases.json}; \texttt{results/T15/initial\_distribution\_sweep.csv} & TBL-2:repo\_asset; Table 2 ([TBL-2]). Arrow / protocol theorem evidence and guardrails: …\\
\texttt{NG\_PROTOCOL\_TRAP} & \texttt{NG\_PROTOCOL\_TRAP.guardrail} & \texttt{results/T15/initial\_distribution\_sweep.csv} & \texttt{results/T15/initial\_distribution\_sweep.csv} & TBL-2:repo\_asset; Table 2 ([TBL-2]). Arrow / protocol theorem evidence and guardrails: …\\
\texttt{NG\_FORCE\_FOREST} & \texttt{NG\_FORCE\_FOREST.guardrail} & \texttt{results/T15/threshold\_sweep.csv} & \texttt{results/T15/threshold\_sweep.csv}; \texttt{results/T14/primitive\_coverage.json} & TBL-3:repo\_asset; Table 3 ([TBL-3]). Graph / affinity theorem evidence and guardrails: …\\
\texttt{NG\_FORCE\_NULL} & \texttt{NG\_FORCE\_NULL.guardrail} & \texttt{results/T15/threshold\_sweep.csv} & \texttt{results/T15/threshold\_sweep.csv}; \texttt{results/T14/primitive\_coverage.json} & TBL-3:repo\_asset; Table 3 ([TBL-3]). Graph / affinity theorem evidence and guardrails: …\\
\texttt{NG\_MACRO\_CLOSURE\_DEFICIT} & \texttt{NG\_MACRO\_CLOSURE\_DEFICIT.guardrail} & \texttt{results/T09/grid\_check.json} & \texttt{results/T09/grid\_check.json}; \texttt{results/T09/closure\_metrics.json} & TBL-4:repo\_asset; Table 4 ([TBL-4]). Closure theorem evidence and provenance: the direc…\\
\texttt{NG\_OBJECT\_CONTRACTIVE} & \texttt{NG\_OBJECT\_CONTRACTIVE.guardrail} & \texttt{results/T15/objecthood\_metric\_sweep.csv} & \texttt{results/T15/objecthood\_metric\_sweep.csv}; \texttt{results/T10/stability\_grid.json} & TBL-5:repo\_asset; Table 5 ([TBL-5]). Objecthood admissible regime and guardrail matrix:…\\
\texttt{NG\_LADDER\_BOUNDED\_INTERFACE} & \texttt{NG\_LADDER\_BOUNDED\_INTERFACE.guardrail} & \texttt{results/T11/no\_ladder\_demos.json} & \texttt{results/T11/no\_ladder\_demos.json}; \texttt{results/T11/definability\_metrics.json} & FIG-6:repo\_asset; Figure 6 ([FIG-6]). Ladder front: fixed idempotent packaging saturate…\\
\end{longtable}
}
 
{\footnotesize
\begin{longtable}{>{\raggedright\arraybackslash}p{0.10\textwidth} >{\raggedright\arraybackslash}p{0.18\textwidth} >{\raggedright\arraybackslash}p{0.20\textwidth} >{\raggedright\arraybackslash}p{0.22\textwidth} >{\raggedright\arraybackslash}p{0.20\textwidth}}
\caption{Frozen callout bindings and caption decisions.}\label{tab:callout-map}\\
\toprule
Callout & Theorem ids & Primary source & Supporting refs & Caption / note\\
\midrule
\endfirsthead
\toprule
Callout & Theorem ids & Primary source & Supporting refs & Caption / note\\
\midrule
\endhead
\midrule
\multicolumn{5}{r}{Continued on next page}\\
\midrule
\endfoot
\bottomrule
\endlastfoot
\texttt{FIG-1} & none & \texttt{authored:theory-package-and-audit-convention-schematic} & \texttt{docs/project/theorem\_atlas.yaml}; \texttt{docs/project/claim\_audit\_freeze.yaml} & Figure 1 ([FIG-1]). Finite audited theory package and audit-convention schematic: microstate s…\\
\texttt{FIG-2} & \texttt{NG\_ARROW\_DPI}; \texttt{NG\_PROTOCOL\_TRAP} & \texttt{results/T15/horizon\_lens\_sweep.csv} & \texttt{results/T08/arrow\_metrics.json}; \texttt{results/T15/initial\_distribution\_sweep.csv} & Figure 2 ([FIG-2]). Hidden-clock DPI / protocol matrix: honest observed path-reversal KL can s…\\
\texttt{FIG-3} & \texttt{NG\_FORCE\_FOREST}; \texttt{NG\_FORCE\_NULL} & \texttt{results/T14/primitive\_coverage.json} & \texttt{results/master/witness\_metrics.csv} & Figure 3 ([FIG-3]). Graph / affinity contrasts: forest support removes cycle-space obstruction…\\
\texttt{FIG-4} & \texttt{NG\_MACRO\_CLOSURE\_DEFICIT} & \texttt{results/T09/closure\_metrics.json} & \texttt{results/T09/grid\_check.json} & Figure 4 ([FIG-4]). Closure front: exact packaged future laws define the finite variational ob…\\
\texttt{FIG-5} & \texttt{NG\_OBJECT\_CONTRACTIVE} & \texttt{results/T10/objecthood\_metrics.json} & \texttt{results/T10/stability\_grid.json} & Figure 5 ([FIG-5]). Objecthood front: in the admissible total-variation regime, contraction yi…\\
\texttt{FIG-6} & \texttt{NG\_LADDER\_IDEM}; \texttt{NG\_LADDER\_BOUNDED\_INTERFACE} & \texttt{results/T11/no\_ladder\_demos.json} & \texttt{results/T11/definability\_metrics.json} & Figure 6 ([FIG-6]). Ladder front: fixed idempotent packaging saturates after one step, bounded…\\
\texttt{TBL-1} & \texttt{NG\_ARROW\_DPI}; \texttt{NG\_PROTOCOL\_TRAP}; \texttt{NG\_FORCE\_FOREST}; \texttt{NG\_FORCE\_NULL}; +4 more & \texttt{results/master/witness\_metrics.csv} & \texttt{results/master/witness\_manifest.json} & Table 1 ([TBL-1]). Theorem front map and support snapshot: the eight theorem fronts, their sec…\\
\texttt{TBL-2} & \texttt{NG\_ARROW\_DPI}; \texttt{NG\_PROTOCOL\_TRAP} & \texttt{results/T15/adversarial\_cases.json} & \texttt{results/T15/initial\_distribution\_sweep.csv} & Table 2 ([TBL-2]). Arrow / protocol theorem evidence and guardrails: the paper binds the arrow…\\
\texttt{TBL-3} & \texttt{NG\_FORCE\_FOREST}; \texttt{NG\_FORCE\_NULL} & \texttt{results/T15/threshold\_sweep.csv} & \texttt{results/T14/primitive\_coverage.json} & Table 3 ([TBL-3]). Graph / affinity theorem evidence and guardrails: exact-support cycle infor…\\
\texttt{TBL-4} & \texttt{NG\_MACRO\_CLOSURE\_DEFICIT} & \texttt{results/T09/grid\_check.json} & \texttt{results/T09/closure\_metrics.json} & Table 4 ([TBL-4]). Closure theorem evidence and provenance: the direct pack uses exact variati…\\
\texttt{TBL-5} & \texttt{NG\_OBJECT\_CONTRACTIVE} & \texttt{results/T15/objecthood\_metric\_sweep.csv} & \texttt{results/T10/stability\_grid.json} & Table 5 ([TBL-5]). Objecthood admissible regime and guardrail matrix: exact contraction metric…\\
\texttt{TBL-6} & \texttt{NG\_LADDER\_IDEM}; \texttt{NG\_LADDER\_BOUNDED\_INTERFACE} & \texttt{results/T11/definability\_metrics.json} & \texttt{results/T11/no\_ladder\_demos.json} & Table 6 ([TBL-6]). Ladder theorem evidence, count core, and corollary map: no-ladder demos, ex…\\
\texttt{TBL-A1} & \texttt{NG\_ARROW\_DPI}; \texttt{NG\_PROTOCOL\_TRAP}; \texttt{NG\_FORCE\_FOREST}; \texttt{NG\_FORCE\_NULL}; +4 more & \texttt{results/master/witness\_manifest.json} & \texttt{results/master/witness\_metrics.csv} & Table A1 ([TBL-A1]). Artifact and evidence provenance map: each frozen core or guardrail claim…\\
\texttt{TBL-A2} & \texttt{NG\_ARROW\_DPI}; \texttt{NG\_PROTOCOL\_TRAP}; \texttt{NG\_FORCE\_FOREST}; \texttt{NG\_FORCE\_NULL}; +4 more & \texttt{derived:lean-support-and-proof-dependency-map} & \texttt{docs/project/readiness\_checklist.yaml}; \texttt{docs/project/theorem\_atlas.yaml} & Table A2 ([TBL-A2]). Lean support and proof dependency map: seven fronts are direct, objecthoo…\\
\end{longtable}
}
  \section{Computational audit and artifact manifest}
The canonical command for rebuilding the computational side of the project is \texttt{make reproduce}.
That command drives the reproducibility pipeline, refreshes theorem-pack summaries, Lean summaries, and rewrite-lane validations, and writes a manifest at \texttt{results/manifest/artifact\_manifest.json}.
The pipeline is deterministic up to the frozen repo contents and toolchain recorded in the manifest.

{\footnotesize
\begin{longtable}{>{\raggedright\arraybackslash}p{0.30\textwidth} >{\raggedright\arraybackslash}p{0.58\textwidth}}
\caption{Reproducibility manifest overview.}\label{tab:manifest-overview}\\
\toprule
Field & Value\\
\midrule
\endfirsthead
\toprule
Field & Value\\
\midrule
\endhead
\midrule
\multicolumn{2}{r}{Continued on next page}\\
\midrule
\endfoot
\bottomrule
\endlastfoot
Canonical command & \texttt{make reproduce}\\
Manifest status & success\\
Generated at UTC & \texttt{2026-03-22T07:25:15Z}\\
Total pipeline steps & 62\\
Hashed files & 350\\
Failed steps & 0\\
\end{longtable}
}
 
The table below lists the theorem-pack and Lean-check steps that are directly relevant to theorem evidence in the main paper, together with the summary JSON outputs they refresh.

{\footnotesize
\begin{longtable}{>{\raggedright\arraybackslash}p{0.20\textwidth} >{\raggedright\arraybackslash}p{0.10\textwidth} >{\raggedright\arraybackslash}p{0.32\textwidth} >{\raggedright\arraybackslash}p{0.26\textwidth}}
\caption{Theorem-pack and Lean-check steps relevant to theorem evidence.}\label{tab:relevant-steps}\\
\toprule
Step & Status & Command & Summary outputs\\
\midrule
\endfirsthead
\toprule
Step & Status & Command & Summary outputs\\
\midrule
\endhead
\midrule
\multicolumn{4}{r}{Continued on next page}\\
\midrule
\endfoot
\bottomrule
\endlastfoot
\texttt{run\_t15\_robustness\_suite} & passed & \texttt{python3 scripts/run\_t15\_robustness\_suite.py} & \texttt{results/T15/fragility\_summary.json}\\
\texttt{run\_t22\_dpi\_protocol\_pack} & passed & \texttt{python3 scripts/run\_t22\_dpi\_protocol\_pack.py} & \texttt{results/T22/dpi\_protocol\_pack\_summary.json}\\
\texttt{run\_t23\_graph\_affinity\_pack} & passed & \texttt{python3 scripts/run\_t23\_graph\_affinity\_pack.py} & \texttt{results/T23/graph\_affinity\_pack\_summary.json}\\
\texttt{run\_t24\_closure\_pack} & passed & \texttt{python3 scripts/run\_t24\_closure\_pack.py} & \texttt{results/T24/closure\_pack\_summary.json}\\
\texttt{run\_t25\_objecthood\_pack} & passed & \texttt{python3 scripts/run\_t25\_objecthood\_pack.py} & \texttt{results/T25/objecthood\_pack\_summary.json}\\
\texttt{run\_t26\_bounded\_interface\_pack} & passed & \texttt{python3 scripts/run\_t26\_bounded\_interface\_pack.py} & \texttt{results/T26/bounded\_interface\_pack\_summary.json}\\
\texttt{run\_t27\_lean\_graph\_bridge} & passed & \texttt{python3 scripts/run\_t27\_lean\_graph\_bridge.py} & \texttt{results/T27/lean\_graph\_bridge\_summary.json}\\
\texttt{run\_t28\_lean\_finite\_lens\_count} & passed & \texttt{python3 scripts/run\_t28\_lean\_finite\_lens\_count.py} & \texttt{results/T28/lean\_finite\_lens\_count\_summary.json}\\
\texttt{run\_t29\_lean\_bounded\_interface\_corollary} & passed & \texttt{python3 scripts/run\_t29\_lean\_bounded\_interface\_corollary.py} & \texttt{results/T29/lean\_bounded\_interface\_corollary\_summary.json}\\
\texttt{run\_t30\_lean\_objecthood\_uniqueness} & passed & \texttt{python3 scripts/run\_t30\_lean\_objecthood\_uniqueness.py} & \texttt{results/T30/lean\_objecthood\_uniqueness\_summary.json}\\
\texttt{run\_t31\_paper\_skeleton} & passed & \texttt{python3 scripts/run\_t31\_paper\_skeleton.py} & \texttt{results/T31/paper\_skeleton\_summary.json}\\
\texttt{run\_t32\_asset\_freeze} & passed & \texttt{python3 scripts/run\_t32\_asset\_freeze.py} & \texttt{results/T32/asset\_freeze\_summary.json}\\
\texttt{run\_t33\_draft\_v1} & passed & \texttt{python3 scripts/run\_t33\_draft\_v1.py} & \texttt{results/T33/draft\_v1\_summary.json}\\
\texttt{run\_t34\_red\_team\_review} & passed & \texttt{python3 scripts/run\_t34\_red\_team\_review.py} & \texttt{results/T34/red\_team\_summary.json}\\
\texttt{run\_t35\_scope\_charter} & passed & \texttt{python3 scripts/run\_t35\_scope\_charter.py} & \texttt{results/T35/lean\_probability\_scope\_charter\_summary.json}\\
\texttt{run\_t36\_lean\_probability\_core} & passed & \texttt{python3 scripts/run\_t36\_lean\_probability\_core.py} & \texttt{results/T36/lean\_probability\_core\_summary.json}\\
\texttt{run\_t37\_lean\_kl\_dpi\_core} & passed & \texttt{python3 scripts/run\_t37\_lean\_kl\_dpi\_core.py} & \texttt{results/T37/lean\_kl\_dpi\_core\_summary.json}\\
\texttt{run\_t38\_lean\_arrow\_dpi} & passed & \texttt{python3 scripts/run\_t38\_lean\_arrow\_dpi.py} & \texttt{results/T38/lean\_arrow\_dpi\_summary.json}\\
\texttt{run\_t39\_lean\_protocol\_trap} & passed & \texttt{python3 scripts/run\_t39\_lean\_protocol\_trap.py} & \texttt{results/T39/lean\_protocol\_trap\_summary.json}\\
\texttt{run\_t40\_closure\_feasibility} & passed & \texttt{python3 scripts/run\_t40\_closure\_feasibility.py} & \texttt{results/T40/lean\_closure\_feasibility\_summary.json}\\
\texttt{run\_t41\_lean\_closure\_variational\_core} & passed & \texttt{python3 scripts/run\_t41\_lean\_closure\_variational\_core.py} & \texttt{results/T41/lean\_closure\_variational\_core\_summary.json}\\
\texttt{run\_t42\_lean\_closure\_direct\_pack} & passed & \texttt{python3 scripts/run\_t42\_lean\_closure\_direct\_pack.py} & \texttt{results/T42/lean\_closure\_direct\_pack\_summary.json}\\
\texttt{run\_t43\_objecthood\_tv\_feasibility} & passed & \texttt{python3 scripts/run\_t43\_objecthood\_tv\_feasibility.py} & \texttt{results/T43/lean\_objecthood\_tv\_feasibility\_summary.json}\\
\texttt{run\_t44\_lean\_objecthood\_direct\_pack} & passed & \texttt{python3 scripts/run\_t44\_lean\_objecthood\_direct\_pack.py} & \texttt{results/T44/lean\_objecthood\_direct\_pack\_summary.json}\\
\end{longtable}
}
  \section{Fragility diagnostics}
The main paper states theorem objects in exact finite terms.
The fragility layer records where nearby proxy choices can distort interpretation without changing the underlying theorem statements.
The frozen source for this section is \texttt{results/T15/fragility\_summary.json}, together with the guardrail claims in \texttt{docs/project/claim\_audit\_freeze.yaml}.

The dominant risks are: changing the initial law in stationary arrow claims, changing the observation lens, fitting first-order macro proxies, thresholding exact support graphs, and switching to non-admissible objecthood metrics.
Each such move is a scope change or a proxy diagnostic, not an in-theorem modification of the audited object.

{\footnotesize
\begin{longtable}{>{\raggedright\arraybackslash}p{0.20\textwidth} >{\raggedright\arraybackslash}p{0.08\textwidth} >{\raggedright\arraybackslash}p{0.20\textwidth} >{\raggedright\arraybackslash}p{0.18\textwidth} >{\raggedright\arraybackslash}p{0.22\textwidth}}
\caption{Fragility items and their frozen guardrails.}\label{tab:fragility-map}\\
\toprule
Fragility item & Severity & Affected claims & Guardrail claim & Mitigation / note\\
\midrule
\endfirsthead
\toprule
Fragility item & Severity & Affected claims & Guardrail claim & Mitigation / note\\
\midrule
\endhead
\midrule
\multicolumn{5}{r}{Continued on next page}\\
\midrule
\endfoot
\bottomrule
\endlastfoot
\texttt{stationary\_zero\_arrow\_is\_fragile\_to\_initial\_distribution} & high & \texttt{NG\_PROTOCOL\_TRAP.core} & \texttt{NG\_PROTOCOL\_TRAP.guardrail} & Protocol-trap claims must be stated on the lifted chain or honest observed path law, under sta…\\
\texttt{macro\_arrow\_depends\_strongly\_on\_lens\_choice} & high & \texttt{NG\_ARROW\_DPI.core}; \texttt{NG\_PROTOCOL\_TRAP.core} & \texttt{NG\_ARROW\_DPI.guardrail} & Fitted first-order macro proxies are not theorem-level arrow audits and can create or distort …\\
\texttt{fitted\_markov\_proxy\_can\_create\_or\_distort\_macro\_arrow} & high & \texttt{NG\_ARROW\_DPI.guardrail}; \texttt{NG\_PROTOCOL\_TRAP.guardrail}; \texttt{NG\_MACRO\_CLOSURE\_DEFICIT.guardrail} & \texttt{NG\_ARROW\_DPI.guardrail} & Fitted first-order macro proxies are not theorem-level arrow audits and can create or distort …\\
\texttt{threshold\_floors\_can\_introduce\_spurious\_one\_way\_support} & medium & \texttt{NG\_FORCE\_FOREST.guardrail}; \texttt{NG\_FORCE\_NULL.guardrail} & \texttt{NG\_FORCE\_NULL.guardrail} & Thresholding or regularization can distort exactness and affinity interpretation; theorem clai…\\
\texttt{object\_multiplicity\_can\_be\_metric\_artifact} & medium & \texttt{NG\_OBJECT\_CONTRACTIVE.guardrail} & \texttt{NG\_OBJECT\_CONTRACTIVE.guardrail} & Theorem objecthood claims must use admissible total-variation/Dobrushin metric assumptions; pr…\\
\end{longtable}
}
  \section{Formalization support map}
This section summarizes the distribution of Lean support across the eight theorem fronts.
The table is generated from \texttt{docs/project/readiness\_checklist.yaml} together with the frozen Lean bridge files under \texttt{docs/project/lean\_*.yaml}.

{\footnotesize
\begin{longtable}{>{\raggedright\arraybackslash}p{0.18\textwidth} >{\raggedright\arraybackslash}p{0.13\textwidth} >{\raggedright\arraybackslash}p{0.30\textwidth} >{\raggedright\arraybackslash}p{0.27\textwidth}}
\caption{Lean support distribution across theorem fronts.}\label{tab:lean-map}\\
\toprule
Theorem & Lean status & Lean modules & Boundary note\\
\midrule
\endfirsthead
\toprule
Theorem & Lean status & Lean modules & Boundary note\\
\midrule
\endhead
\midrule
\multicolumn{4}{r}{Continued on next page}\\
\midrule
\endfoot
\bottomrule
\endlastfoot
\texttt{NG\_ARROW\_DPI} & direct & \texttt{lean/SixBirdsNoGo/FiniteKLDPI.lean}; \texttt{lean/SixBirdsNoGo/FiniteKLDPIExample.lean}; \texttt{lean/SixBirdsNoGo/ArrowDPI.lean}; +1 more & Direct support recorded in the frozen bridge files.\\
\texttt{NG\_PROTOCOL\_TRAP} & direct & \texttt{lean/SixBirdsNoGo/ArrowDPI.lean}; \texttt{lean/SixBirdsNoGo/ArrowDPIExample.lean}; \texttt{lean/SixBirdsNoGo/ProtocolTrap.lean}; +1 more & Direct support recorded in the frozen bridge files.\\
\texttt{NG\_FORCE\_FOREST} & direct & \texttt{lean/SixBirdsNoGo/ForestForceBridge.lean}; \texttt{lean/SixBirdsNoGo/ForestForceBridgeExample.lean} & Direct support recorded in the frozen bridge files.\\
\texttt{NG\_FORCE\_NULL} & direct & \texttt{lean/SixBirdsNoGo/NullForceBridge.lean}; \texttt{lean/SixBirdsNoGo/NullForceBridgeExample.lean} & Direct support recorded in the frozen bridge files.\\
\texttt{NG\_MACRO\_CLOSURE\_DEFICIT} & direct & \texttt{lean/SixBirdsNoGo/ClosureVariationalCore.lean}; \texttt{lean/SixBirdsNoGo/ClosureVariationalCoreExample.lean}; \texttt{lean/SixBirdsNoGo/ClosureDirectPack.lean}; +1 more & Direct support recorded in the frozen bridge files.\\
\texttt{NG\_OBJECT\_CONTRACTIVE} & auxiliary\_only & \texttt{lean/SixBirdsNoGo/ContractionUniqueness.lean}; \texttt{lean/SixBirdsNoGo/ContractionUniquenessExample.lean} & Boundary note: uniqueness core only; broader TV/Dobrushin plus epsilon-stability package is no…\\
\texttt{NG\_LADDER\_IDEM} & direct & \texttt{lean/SixBirdsNoGo/Idempotence.lean} & Direct support recorded in the frozen bridge files.\\
\texttt{NG\_LADDER\_BOUNDED\_INTERFACE} & direct & \texttt{lean/SixBirdsNoGo/Idempotence.lean}; \texttt{lean/SixBirdsNoGo/FiniteLensDefinability.lean}; \texttt{lean/SixBirdsNoGo/FiniteLensDefinabilityExample.lean}; +2 more & Direct support recorded in the frozen bridge files.\\
\end{longtable}
}
 
For \texttt{NG\_OBJECT\_CONTRACTIVE}, the permanent boundary is explicit.
The uniqueness-from-strict-contraction core is formalized, but the broader total-variation/Dobrushin plus $\varepsilon$-stability package used in the paper is not Lean-direct and remains frozen behind the objecthood TV boundary note.
 \end{document}
